% =========================
% URF Core Insert (Canonical)
% Expander-Exclusion Boundary + Corrected Local→Global Lemma
% =========================

\section{URF Core Boundary: Expander Exclusion and Corrected Local→Global Rigidity}

\subsection{Basic notation}
Let $G=(V,E)$ be a finite, simple graph with maximum degree $\Delta<\infty$.
For $v\in V$ and $r\in\mathbb N$, let
\[
B_r(v):=\{u\in V:\operatorname{dist}(u,v)\le r\},
\qquad
\partial B_r(v):=\{u\in B_r(v):\exists w\notin B_r(v)\text{ with }uw\in E\}.
\]
Define the (edge) Cheeger constant
\[
h(G):=\min_{\emptyset\ne S\subsetneq V,\ |S|\le |V|/2}\frac{|E(S,V\setminus S)|}{|S|}.
\]

\subsection{FO$^k$ local homogeneity (radius-$R$)}
Fix $k\in\mathbb N$ and $R\in\mathbb N$.
Write $\mathrm{tp}^{k}_R(G,v)$ for the FO$^k$ \emph{radius-$R$ local type} of $v$ (Gaifman-normal-form; equivalently, the $k$-variable theory of the induced structure on $B_R(v)$ with distinguished root $v$).

\begin{definition}[FO$^k_R$-local homogeneity]
$G$ is \emph{FO$^k_R$-locally homogeneous} if
\[
\mathrm{tp}^{k}_R(G,u)=\mathrm{tp}^{k}_R(G,v)\quad\forall u,v\in V.
\]
\end{definition}

\subsection{Growth conditions (expander-excluding)}
\begin{definition}[Uniform polynomial growth]
A family $\{G_n\}$ of graphs has \emph{uniform polynomial growth} if there exist constants $C,c>0$ such that
\[
|B_r(v)|\le C\, r^{c}\quad\forall n,\ \forall v\in V(G_n),\ \forall r\le \operatorname{diam}(G_n).
\]
\end{definition}

\begin{definition}[Uniform subexponential growth]
A family $\{G_n\}$ has \emph{uniform subexponential growth} if there exists a function $\eta:\mathbb N\to(0,\infty)$ with $\eta(r)\to 0$ such that
\[
|B_r(v)|\le \exp(\eta(r)\,r)\quad\forall n,\ \forall v\in V(G_n),\ \forall r\le \operatorname{diam}(G_n).
\]
\end{definition}

\subsection{URF Core Axiom: Expander Exclusion (Boundary Principle)}
\begin{axiom}[URF-Core-EX: Growth excludes expansion]
If a bounded-degree graph family $\{G_n\}$ has uniform subexponential growth, then it is not an expander family:
\[
\Delta(G_n)\le \Delta,\ \text{uniform subexp growth}\ \Longrightarrow\ \lim_{n\to\infty} h(G_n)=0.
\]
\end{axiom}

\subsection{Certified base instance (proved): subexponential growth implies small Cheeger constant}
\begin{theorem}[Base theorem: Subexponential growth $\Rightarrow$ no constant expansion]\label{thm:subexp-noexp}
Let $\{G_n\}$ be a family with $\Delta(G_n)\le \Delta$ and uniform subexponential growth.
Then for every $\varepsilon>0$ there exists $n_0$ such that for all $n\ge n_0$,
\[
h(G_n)\le \varepsilon.
\]
In particular, $\{G_n\}$ is not an expander family.
\end{theorem}

\begin{proof}
Fix $\varepsilon>0$.
By uniform subexponential growth, choose $r$ so large that
\[
\exp(\eta(r)\,r) < (1+\varepsilon)^{r}.
\]
Fix $n$ with $\operatorname{diam}(G_n)\ge r$ (otherwise $|V(G_n)|$ is bounded and the claim is trivial).
Pick any $v\in V(G_n)$ and consider the sequence $a_j:=|B_j(v)|$ for $0\le j\le r$.
If $a_{j+1}>(1+\varepsilon)a_j$ for all $0\le j<r$, then
\[
a_r>(1+\varepsilon)^r a_0=(1+\varepsilon)^r,
\]
contradicting $a_r\le \exp(\eta(r)\,r)$.
Hence there exists $0\le j<r$ such that
\[
|B_{j+1}(v)|\le (1+\varepsilon)|B_j(v)|.
\]
Then
\[
|\partial B_j(v)|\le |B_{j+1}(v)\setminus B_j(v)|\le \varepsilon |B_j(v)|.
\]
Let $S:=B_j(v)$. If $|S|\le |V|/2$, then by degree-boundedness,
\[
|E(S,V\setminus S)|\le \Delta\,|\partial S|\le \Delta\,\varepsilon |S|,
\]
so $h(G_n)\le \Delta\,\varepsilon$.
If instead $|S|>|V|/2$, apply the same argument to a ball grown from some vertex outside $S$
(or equivalently take $V\setminus S$ and adjust radii); one obtains a set $S'$ with $|S'|\le |V|/2$
and $|E(S',V\setminus S')|/|S'|\le \Delta\,\varepsilon$.
Since $\varepsilon$ was arbitrary, $h(G_n)\to 0$.
\end{proof}

\subsection{Corrected URF Core Lemma (replacing the false version)}
\begin{lemma}[URF-Core-LG (Corrected): Local homogeneity does not force global simplicity, but forces non-expansion under growth]\label{lem:urf-core-lg-corrected}
Fix $\Delta<\infty$. Let $\{G_n\}$ be a bounded-degree family with uniform subexponential growth.
Then for any fixed $k,R$,
\[
\text{$G_n$ FO$^k_R$-locally homogeneous for all $n$}
\quad\Longrightarrow\quad
\lim_{n\to\infty} h(G_n)=0.
\]
\end{lemma}

\begin{proof}
FO$^k_R$-local homogeneity is not used; the conclusion follows directly from Theorem~\ref{thm:subexp-noexp}.
(Interpretation: the \emph{only} obstruction to the naive Local$\to$Global slogan is the expander regime.)
\end{proof}

\subsection{Redefined global invariant (expander-immune)}
\begin{definition}[Geometric cycle-overlap complexity at scale $R$]
Fix $R\in\mathbb N$.
Let $\mathcal C_{\le 2R}(G)$ be the set of (simple) cycles of length $\le 2R$.
Define
\[
\mathsf{CycRank}_R(G):=\operatorname{rank}_{\mathbb F_2}\Big(\text{cycle-space generated by }\mathcal C_{\le 2R}(G)\Big).
\]
\end{definition}

\paragraph{Expander immunity.}
If $\operatorname{girth}(G)>2R$, then $\mathcal C_{\le 2R}(G)=\emptyset$ and hence
\[
\mathsf{CycRank}_R(G)=0.
\]

\subsection{Counterexample capsule (closure)}
\begin{proposition}[High-girth expanders refute the naive Local$\to$Global lemma]\label{prop:counterexample}
Fix $\Delta\ge 3$ and $R\in\mathbb N$.
There exist $\Delta$-regular graphs $G_n$ with $\operatorname{girth}(G_n)>2R$ and $h(G_n)\ge c(\Delta)>0$.
Moreover, for any fixed $k$,
\[
\mathrm{tp}^{k}_R(G_n,u)=\mathrm{tp}^{k}_R(G_n,v)\quad\forall u,v\in V(G_n),
\]
yet $h(G_n)$ is bounded below.
\end{proposition}
